\section{Introduction}\label{sec:intro}
A quadplane carries a fixed wing with a pusher propeller and four lift rotors on booms. Between hover and cruise it must accelerate through the airspeed range in which the wing produces little lift, hand the weight of the aircraft from the rotors to the wing and shut the rotors down, and reverse the sequence to land. This hand-off is the phase in which small vertical take-off and landing (VTOL) aircraft are most often lost, and the phase that open-source autopilots handle with the least sophistication: the PX4 stack, which flew the aircraft studied here, ramps the pusher and cuts the rotors when a one-hertz airspeed sample crosses a threshold~\cite{Meier2015,PX4docs}. The threshold says nothing about the lift actually carried by wing and rotors together, which is the quantity that decides whether the aircraft flies or falls, and it is tuned by hand for one aircraft in one condition.

Learning-based controllers have reached the flight of quadrotors and fixed-wing aircraft~\cite{Kaufmann2023,Song2023}, and neural networks that imitate or replace model predictive controllers are established~\cite{Hertneck2018,Karg2020,Chen2018}. Two obstacles keep them away from the hand-off. First, the plant is uncertain: the polar of a foam wing, the gain and lag of the attitude loop and the maps and lags of two thrust actuators differ from one airframe to the next and from one flight to the next, so a policy trained on one model has no reason to work on another. Second, the hand-off is safety-critical, and a network offers no guarantee on the lift margin, the angle of attack or the sink rate. The safe-learning literature answers the second obstacle with safety filters built on control barrier functions~\cite{Ames2019}, predictive safety filters~\cite{Wabersich2021} and backup controllers~\cite{Gurriet2020}, and the first with domain randomisation~\cite{Tobin2017,Peng2018}; what is missing is a controller that combines both for a plant family with lagged actuators and an attitude loop inside the plant, and that is validated on a flight-identified aircraft against the full range of alternatives.

\emph{Novelty.} This article treats the hand-off as a learning problem across a family of plants: a single policy must fly every aircraft of an identified parameter box without being told which one it is flying. What is new is (a) a history-conditioned transformer, HOT, whose context token is trained to carry the plant it infers from the last $L$ autopilot samples, so that the policy adapts within a fraction of a second without online identification; (b) a lift-safety module with a proof of recursive feasibility: a backup manoeuvre built from measured quantities only is rolled out from the predicted next state over the vertices of the box and the plants consistent with the observed transitions, the network command being accepted only if every rollout keeps the lift margin, the angle of attack, the flight path and the sink rate within bounds and settles into level flight; and (c) physics-informed training through the differentiable identified model under domain randomisation over the box, gusts and sensor noise, with no expert controller, no reward engineering and no simulator beyond the identified equations. The contributions are: (i) the HOT architecture with its four modules and a bounded head that enforces the input box and the actuator rates by construction (Section~\ref{sec:method}); (ii) the lift-safety module and its analysis, Lemmas~\ref{lem:corner}--\ref{lem:cover} and Theorem~\ref{thm:main}, with a verifiable convergence condition and an identifiability result for the context token (Section~\ref{sec:analysis}); (iii) a flight-identified platform, an \SI{8}{kg} quadplane whose ten-parameter box comes from its PX4 log (Section~\ref{sec:platform}); and (iv) an evaluation against nine controllers, from the autopilot rule to online model predictive control, feed-forward and recurrent policies and two reinforcement-learning agents, on the nominal plant, on \numCross{} unseen plants of the box, outside the box, under gusts and noise, from the states recorded in flight and with the PX4 firmware in the loop, with Wilcoxon and Friedman statistics, ablations, parameter sweeps and attention analyses (Section~\ref{sec:experiments}). The flight data are released as a dataset~\cite{Dony2026data}. The code, the trained networks, the results of every experiment and the article sources are available at \url{https://github.com/AdamDony/quadplane-hot}.

\emph{Outline.} Section~\ref{sec:related} places the article among transition controllers and learned controllers with safety. Section~\ref{sec:prelim} fixes the notation, the longitudinal model and the facts used later. Section~\ref{sec:problem} states the plant family, the constraints and the objectives. Section~\ref{sec:method} presents the four modules of HOT, the lift-safety module, the training procedure and the two algorithms. Section~\ref{sec:analysis} proves the lemmas, the main theorem and the identifiability result, and lists the hypotheses verified numerically. Section~\ref{sec:platform} describes the aircraft, its identified box, the simulator and the firmware-in-the-loop set-up. Section~\ref{sec:experiments} reports the experiments, statistics, ablations, interpretability, time cost and the comparison with previous work, and Section~\ref{sec:conclusion} concludes.
